\section{Synchronization}

\begin{example}
    A famous example of thread-safety issues is the counter. When a counter c is incremented using c++ or ++c the value of c
    first gets read, then incremented and then stored. The counter fails if one thread stores his result between
    the time another reads and stores again.
\end{example}

\begin{definition}
    A \textbf{Bad Interleaving} is an interleaving of operations from multiple threads that leads to an incorrect
    program state.
\end{definition}

\begin{definition}
    A \textbf{Critical Section} is a codeblock that might lead to bad interleavings. To ensure program correctness we need the enforce mutual exclusion on them.
\end{definition}

\begin{definition}
    \textbf{Thread safty} refers to nothing bad happeing in all possible interleavings and thus implies program safety.
\end{definition}

\begin{definition}
    \textbf{Thread Liveness} refers to something good happening eventually, which implies the correctness of a program.
\end{definition}

\subsection{Race Conditions}

\begin{definition}
    A \textbf{Race condition} occurs in concurent programming when the correctness of the system depends on
    the specific interleaving or ordering of execution of operations.
\end{definition}

\begin{definition}
    A \textbf{Low-level race condition} is a race condition that occurs because of insufficient
    synchronization of memory access (Data Races).
\end{definition}

\begin{definition}
    A \textbf{High-level race condition} is a race condition that
    occurs because of a wrong decomposition/algorithm. The program might ensure mutual exclusion on the
    shared resources but the program might still produce wrong results due to \textbf{bad interleavings}.
\end{definition}

Generally we want one of three things when it comes to shared ressources:
\begin{itemize}
    \item \textbf{Thread local (Isolated Mutability):} We dont use the ressource of other threads.
    \item \textbf{Read only (Immutable):} We only read and do not change the ressource
    \item \textbf{Mutual exclusion (Synchronization):} We use the ressource of other threads but we ensure that only one thread accesses it at a time.
\end{itemize}


\subsection{Locks and Synchronized}

Every Java-Object has an \textbf{intrinsic lock or lock monitor}. These locks are unique and can be optaind by
threads which allows us to lock a certain codeblock. If a thread wants to access such a codeblock and
execute the code within it first has to optain the lock, while it has the lock, no other thread can accuire
it.

\begin{codeexample}
synchronized(lockObject) {
    i++;
}
\end{codeexample}

This ensures mutual exlusion on the synchronized codeblock. It is also possible to use synchronized as a
keyword in a function to sync. the whole function.

\begin{remark}
    By default synchronized blocks use "this" as the lock-object and unlock when an exception occures.
\end{remark}

\begin{remark}
    Java-Locks are \textbf{reentrant}, meaning they can be locked multiple times by the same thread recursively (it
    will then have to release the lock as many times as it acquired it).
\end{remark}

For incrementing and decrementing integers there is a special DataType AtomicInteger which has built in
synchronization as incrementing/decrementing is done in a single step.

\begin{codeexample}
AtomicInteger counter = new AtomicInteger(0);
counter.incrementAndGet();
\end{codeexample}

\subsection{Wait, notify, notifyAll}

We can use wait() on a lock-object to tell the current thread to sleep until further notice if the lock is
not available. This allows use to save ressources.

\begin{codeexample}
synchronized(lockObject) {
    while(!condition) {
        lockObject.wait();
    }
    // do work  
}
\end{codeexample}

To wake one or multiple threads and tell them that the lock is now available we can use notify() or
notifyAll() on the lock-object to wake up an aribtrary waiting thread or all of them.

\begin{important}
    Always \textbf{wait inside a while loop}, as this ensures checking the condition after beeing notified.
\end{important}

\begin{codeexample}
synchronized(lockObject) {
    while (!condition) {
        lockObject.wait();
    }
    //do work
    notifyAll();
}       
\end{codeexample}

\subsection{Lock API}

\begin{example}
    The problem with the below code is that it is \textbf{not thread-safe}. If two threads try to withdraw money
    at the same time, they might both read the same balance and then both write to it, resulting in a
    wrong or even negative final balance. We need mutual exclusion in that critical section.
\end{example}

\begin{codeexample}
class BankAccount {
    private int balance;

    int balance() {return balance;}

    void setBalance (in x) {balance = x;}

    void withdraw(int amount) {
        int b = balance();
        if (b less than amount) {
            throw new Exception("Not enough money");
        }
        setBalance(b-amount);
    }
}
\end{codeexample}

Instead of using built-in synchronized blocks we can also use locks from the java.util.concurrenct.locks package.

\begin{codeexample}
class BankAccount {
    private Lock lock = new ReentrantLock();
    ...
    void withdraw(int amount) {
        lock.lock();
        try {
            ...
        } finally {
            lock.unlock();
        }
    }
}
            
\end{codeexample}

\begin{important}
    Always use a try-finally block to ensure the lock is released even if an error occures.            
\end{important}

\begin{remark}
    There is also a \textbf{lock.tryLock(100, TimeUnit.MILLISECONDS)} method which tries to accuire the lock only for
    a given amount of time. This is useful for avoiding deadlocks (ex. when two threads have each one lock
    and wait on the other to get released)
\end{remark}

\begin{tabular}{|c|c|c|}
\hline
\textbf{Aspect} & \textbf{synchronized} & \textbf{Lock API} \\\hline
Lock handling & Automatic (JVM managed) & Manual (lock() / unlock()) \\\hline
Lock release & Automatic & Must be done explicitly (finally!) \\\hline
Scope & Method or block scoped & Flexible (can span multiple methods) \\\hline
Waiting behavior & Thread blocks indefinitely & tryLock() avoids blocking \\\hline
Interruptibility & Not interruptible while waiting & lockInterruptibly() supports interrupt \\\hline
Flexibility & Low & High \\\hline
Complexity & Simple, less error-prone & More complex, error-prone \\\hline
Deadlock risk & Lower & Higher (if unlock forgotten) \\\hline
Typical use & Simple synchronization & Advanced concurrency control \\\hline
\end{tabular}



\subsection{Guidelines for Concurrent programming}

\begin{important}
    \textbf{Thread-Locality:} Minimize the sharing of mutable state between threads. Only share ressources if needed.
\end{important}

\begin{important}
    \textbf{Immutability:} Use immutable data structures whenever possible. They are inherently thread-safe.
\end{important}

\begin{important}
    \textbf{Consitent locking:} If you need to share mutable state, use locks to ensure mutual exclusion.
\end{important}

\begin{definition}
    \textbf{Lock granularity } refers to the amount of code that is proteted by a lock in ways:
    \begin{itemize}
        \item Size of synchronized blocks
        \item Number of objects a lock is used for
    \end{itemize}
\end{definition}

\begin{theorem}
    \textbf{Coarse-grained locking} is easier to implement but leads to less concurrency and therfore performance loss. \textbf{Fine-grained locking} is harder to implement without creating bugs but leads to more concurrency and therfore performance gain.
\end{theorem}

\begin{important}
    \textbf{Start with coarse-grained (simpler)} and move to fine-grained (performance) only if contention on the coarser locks becomes an issue.
\end{important}

\begin{important}
    Do not do expensive computations or I/O in critical sections, but also don’t introduce race conditions
\end{important}

\begin{definition}
    \textbf{Atomicity: } An operation is atomic if it is performed as a single, indivisible unit. In other words, it either completes entirely or not at all, without any intermediate states being visible to other threads.
\end{definition}
